%% kapa: exam=1 round=26 year=2015 period=2 form=A subject=appraisal_law subject_ko=부동산관계법규 seq=18 number=18
%% source: Q-Net 감정평가사 기출문제 (cst00309) artlSeq=5138001 file=extracted/1008/5138001/제26회_감정평가사 1차_기출/제26회_감정평가사 1차_기출/2교시_감평_A형.pdf
%% conversion: kordoc 4.24.51 -> tools/gen1.py (auto); review=proofread-sonnet; flags=none
\begin{question}{18}
부동산가격공시 및 감정평가에 관한 법령상 감정평가사징계위원회(이하 ‘징계위원회’라 함)에 관한 설명으로 옳은 것은?
\begin{choices}
  \item 국토교통부의 5급 공무원은 징계위원회의 위원이 될 수 없다.
  \item 「고등교육법」에 따른 대학에서 토지ㆍ주택 등에 관한 이론을 가르치는 조교수 이상은 징계위원회의 위원이 될 수 없다.
  \item 징계위원회의 위원의 임기는 2년으로 하되, 연임할 수 없다.
  \item 징계위원회는 9명의 위원으로 구성하며, 징계위원회의 의결은 재적위원 과반수의 찬성으로 한다.
  \item 징계위원회의 부위원장은 위원장이 지명하는 자로 한다.
\end{choices}
\end{question}
\answer{1}
